\section{Commutativity theorem}
In this section, we provide a commutativity criterion for formal group laws. More precisely,

\begin{theorem}[{\cite[Chap.~1 Theorem 6.1]{hazewinkel1978formal}}]
    Let $R$ be a ring. Then every one-dimensional formal group law over $R$ is commutative if
    and only if $R$ has no nonzero element $a$ that is both nilpotent and
    $\ZZ$-torsion.
    \label{comm_thm}
\end{theorem}

To prove this theorem, we need some auxiliary lemmas.

\begin{lemma}
    Let $F$ be a formal group law over a commutative ring $R$, then we have the following results:
    \begin{enumerate}
        \item $F$ is commutative if and only if $F(X, F(Y, i(X))) = Y$.
        \item $F$ is commutative if and only if $F(X, F(Y, F(i(X), i(Y)))) = 0$.
    \end{enumerate}
    \label{lem: comm_iff}
\end{lemma}

\begin{proof}
    By definition,
    \[
    \begin{aligned}
        F \text{is commutative} \iff F(X,Y) = F(Y,X).
    \end{aligned} \]
    \textit{Proof of Part (1):}
    \[
    \begin{aligned}
        F(X,Y) = F(Y,X) & \iff F(F(X,Y), i(X)) = F(F(Y,X), i(X)) \\
        & \iff F(X, F(Y, i(X))) = Y.
    \end{aligned}\]

    \textit{Proof of Part (2):}
    \[
    \begin{aligned}
        F(X,Y) = F(Y,X) & \iff F(X, F(Y, i(X))) = Y \\
        & \iff F(F(X, F(Y, i(X))), i(Y)) = F(Y, i(Y)) \\
        & \iff F(X, F(Y, F(i(X), i(Y)))) = 0.
    \end{aligned}\]
\end{proof}

\begin{lemma}[{\cite[Chap.~1 Lemma 6.2.1]{hazewinkel1978formal}}]
    Let $F(X,Y)$ be a non-commutative one-dimensional formal group law over a ring
    $R$. Then there exists a nonzero homomorphism from $F(X,Y)$ into $\GG_a(X,Y)$ or into
    $\GG_m(X,Y)$.
    \label{lem: non_com_hom}
\end{lemma}


\begin{proof}
    For simplicity, we denote $F(X,Y)$ by $X +_F Y$ and its inverse $i(X)$ by $-_F X$.
    Let $H(X,Y) = F(X, F(Y, i(X)))$.
    By Lemma \ref{lem: comm_iff} and $F(X,Y)$ is non-commutative,
    we have
    \[H(X,Y) = Y + \sum _{i=1}^\infty r_n(X) Y ^ n,\]
    with $r_n(X) \neq 0$ for some $n$. Since $H(0,Y) = F(0, F(Y, i(0))) = Y$, we have $r_n(0) = 0$ for all $n$. By
    associativity of $F(X,Y)$, we have
    \[
    \begin{aligned}
        H(X_1,H(X_2,Y)) &= X_1 +_F H(X_2,Y) +_F (-_F (X_1)) \\
        & = X_1 +_F (X_2 +_F Y +_F (-_F X_2)) +_F (-_F (X_1)) \\
        & = (X_1 +_F X_2) +_F Y +_F (-_F (X_1 +_F X_2)) \\
        &= H(X_1+_F X_2, Y)
    \end{aligned}
    \]
    % \begin{equation}
    %     H(X_1,H(X_2,Y)) = H(X_1+_F X_2, Y)
    %     \label{eq: H eq}
    % \end{equation}

    Let $i$ be the least integer such that $r_i(X) \neq 0$.

    \textit{For the case $i = 1$:} we have
    \begin{equation*}
        H(X,Y) \equiv (1+ r_1(X)) Y \pmod{Y ^ 2}
    \end{equation*}

    Combining this with equation $H(X_1,H(X_2,Y)) = H(X_1+_F X_2, Y)$, we have
    \[
    \begin{aligned}
        H(X_1, H(X_2,Y)) &\equiv (1 + r_1(X_1)) H(X_2,Y) \\
         &\equiv (1+ r_1(X_1)) (1 + r_1(X_2)) Y
        \pmod{Y ^ 2}, \\
        H(X_1+_F X_2, Y) &\equiv (1+ r_1(X_1+_F X_2)) Y \pmod{Y ^ 2}.
    \end{aligned}\]

    Therefore, $(1 + r_1(X_1)) (1+ r_1(X_2)) = 1 + r_1(X_1 + _F X_2)$, i.e.
    \[r_1(X_1+_F X_2) = r_1(X_1) + r_1(X_2) + r_1(X_1) r_1(X_2).\]

    This means $r_1(X)$ is a formal group homomorphism from $F$ to $\GG_m$.

    \textit{For the case $i > 1$:} we have

    \[H(X,Y) \equiv Y + r_i(X) Y ^ i \pmod{Y ^ {i+ 1}}.\]

    Note that we have
    \[
    \begin{aligned}
    H(X_1, H(X_2,Y)) &\equiv H(X_2,Y) + r_i(X_1) H(X_2,Y)^i \\
        &\equiv Y + r_i(X_2)Y^i + r_i(X_1)Y^i \pmod{Y^{i+1}}, \\
    H(X_1 +_F X_2,Y) &\equiv Y + r_i(X_1 +_F X_2)Y^i \pmod{Y^{i+1}}, \\
    H(X_1,H(X_2,Y)) &= H(X_1 +_F X_2,Y).
    \end{aligned}
    \]

    Therefore, $r_i(X_1+_F X_2) = r_i(X_1) + r_i(X_2)$, which means $r_i(X)$ is a formal group
    homomorphism from $F$ to $\GG_a$.
\end{proof}

\begin{lemma}[{\cite[Chap.~1 Lemma 6.2.11]{hazewinkel1978formal}}]
    Let $R$ be an integral domain, and let $F, F'$ be one-dimensional formal group laws
    over $R$. Suppose there exists a homomorphism
    $\alpha(X): F(X,Y) \to F'(X,Y)$ and $F'(X,Y)$ is commutative. Then $F(X,Y)$ is
    commutative.
    \label{lem: target_com}
\end{lemma}

\begin{proof}
    Let $r$ be the order of $\alpha (X) = \sum_{n=1} ^ \infty a_n X^n$, i.e. the least $n$ such
    that $a_n$ is nonzero. Let $C(X,Y) = X +_F Y +_F (-_F X) +_F (-_F Y) $ be the
    commutator of $X$ and $Y$ with respect to $F$. \par
    By Lemma \ref{lem: comm_iff}, a formal group law $F$ is commutative if and only if the commutator
    $C(X,Y)$ with respect to $F$ is zero.

    Since $F'$ is commutative, we have
    \[ \alpha (C(X,Y)) = \alpha(X)+_G \alpha(Y)+_G(-_G \alpha(X))+_G(-_G\alpha(Y)) = 0.\]
    Suppose $F$ is non-commutative, i.e. $C(X,Y) \neq 0$. Let $C_n(X,Y)$ be the homogeneous
    component in $C(X,Y)$ of degree $n$, and let $m$ be the least integer such that
    $C_m(X,Y)$ is non-zero. We have
    \[ \alpha (C(X,Y)) \equiv a_r C_m(X,Y) ^ r \pmod{\deg (mr + 1)}. \]
    Since $a_r \in R$ and $R$ is an integral domain, this term is non-zero, which contradicts
    the assumption that $\alpha (C(X,Y))= 0$.
\end{proof}

% \begin{lemma}[Necessity of the condition on $R$]
%     Let $0 \neq a \in R$ be an element in $R$ is a nilpotent and $\ZZ$-torsion, then
%     there is formal group law $F(X,Y)$ over $R$ such that $F(X,Y)$ is
%     non-commutative.
% \end{lemma}

% here Additive torsion free is equivalent to characteristic zero.
% \begin{lemma}[Sufficiency in case $R$ is a $\ZZ$-torsion free]
%     Suppose $R$ is $\ZZ$-torsion free, then every one dimensional formal group law over ring $R$ is
%     commutative.
% \end{lemma}

% \begin{lemma}[Sufficiency in case $R$ is integral domain]
%     Let $R$ be a integral domain, then every one dimensional formal group law over ring $R$ is
%     commutative.
% \end{lemma}

% \begin{proof}
%     Let $F(X,Y)$ be a formal group law over $R$, suppose that $F(X,Y)$ is non-commutative, then by
%     lemma $\ref{lem: non_com_hom}$, there is a non zero formal group homomorphism from $F(X,Y)$ to
%     $\GG_a(X,Y)$ or $\GG_m(X,Y)$. According to lemma $\ref{lem: target_com}$ and
%     $\GG_a(X,Y), \GG_m (X,Y)$ is commutative formal group law, then $F(X,Y)$ is commutative. This
%     leads to a contradiction.
% \end{proof}

\begin{proof}[Proof of Theorem \ref{comm_thm}]
    First, we prove the necessity of the condition on $R$.\par
    Let $a \in R$ be an nonzero element of $R$ that is nilpotent and $\ZZ$-torsion. We need to show
    that there exists
    a formal group law $F(X,Y)$ over $R$ such that $F(X,Y)$ is non-commutative. \par
    Let $n$ be the least positive integer such that $n a = 0$. Then $n>1$ and there exists a prime number
    $p$ such that $p \mid n$. Let $b = \frac{n}{p} a$. Then $b\neq 0$ and $p b = 0$, since $n$ is
    minimal. Let $m$ be the least integer such that $b^m= 0$ and let $c = b^{m-1}$. Therefore, this $c$
    satisfies $c \neq 0, c ^ 2 = 0, p c = 0$. Denote $F(X,Y) = X + Y + c X Y ^ p$. Then
    \begin{equation*}
        F(F(X,Y), Z) = X + Y + Z + c(X Y^p+ XZ^p + YZ^p) = F(X,F(Y,Z)).
    \end{equation*}
    Therefore, $F(X,Y)$ is a non-commutative formal group law over $R$.

    Now we prove sufficiency of the condition on $R$. Suppose first that $R$ is a integral domain.
    Let $F(X,Y)$ be a formal
    group law over $R$, where $R$ is an integral domain. Suppose that $F(X,Y)$ is non-commutative.
    Then, by Lemma $\ref{lem: non_com_hom}$, there exists a nonzero formal group homomorphism from $F(X,Y)$ to
    $\GG_a(X,Y)$ or $\GG_m(X,Y)$. Since $\GG_a(X,Y)$ and $\GG_m (X,Y)$ are commutative formal group laws,
    Lemma $\ref{lem: target_com}$ implies that $F(X,Y)$ is commutative. This
    leads to a contradiction.\par

    Suppose that $R$ is $\ZZ$-torsion free. A proof that $F$ is commutative over $R$ is given
    later in Corollary \ref{cor: torsion_free_comm}.
    Also, there is a different proof in {\cite[Chap.~1 Lemma 6.1.3]{hazewinkel1978formal}}. \par

    Finally, let $R$ be a ring such that $R$ has no nonzero element $a$ that is both nilpotent and
    $\ZZ$-torsion.

    Given a prime ideal $\gothp$ in $R$, consider the ring homomorphism $\phi : R \to R / \gothp$.
    Since $R / \gothp$ is an integral domain, $\phi_* F(X,Y)$ is commutative. If we denote
    $F(X,Y) = \sum_{i, j} a_{ij} X ^ i Y ^ j$, then it follows that
    $\phi_* (a_{ij} - a_{ji}) = 0 \in R / \gothp$, which is equivalent to
    $a_{ij} - a_{ji} \in \gothp$. Since this holds for all prime ideals $\gothp$, we have
    \begin{equation*}
        a_{ij} - a_{ji} \in \bigcap _{\gothp \ prime} \gothp.
    \end{equation*}

    Therefore, for any $i, j$, $a_{ij} - a_{ji}$ is nilpotent in $R$.

    Consider the ring homomorphism $\phi_\infty : R \to R \otimes_{\ZZ} \QQ$. Since
    $R \otimes_\ZZ \QQ$ is $\ZZ$-torsion free, $\phi_{\infty *} F(X,Y)$ is commutative.
    Therefore, $\phi_\infty (a_{ij} - a_{ji}) = 0 \in R \otimes_\ZZ \QQ$, which means
    $a_{ij} -a_{ji}$ is $\ZZ$-torsion. Since there is no nonzero element in $R$ that is
    nilpotent and $\ZZ$-torsion at the same time, it follows that for any $i, j$,
    $a_{ij} - a_{ji} = 0$, which implies that $F(X,Y)$ is commutative.

\end{proof}